% figure:
\begin{figure}[h] %h : ottimizza lo spazio, H : mette figura nel punto specificato nel codice
\caption{Esempio di switch su connessioni TX ed RX del modulo ESP8266}
\centering
\includegraphics[width=\textwidth]{figures/SwitchEs.png}
\end{figure}
%-----------------------------------------------------------------------

%riferimento ad una parte di testo o label
%colorato blu:
\riferimento{riferimento}{testo cliccabile}
%colorato a piacere
\riferimento[colore]{riferimento}{testo cliccabile}
%normale
\hyperref[riferiento]{testo cliccabile}
%-----------------------------------------------------------------------

%tabella
\begin{center} \begin{tabular}{|c l|} %c = center, l = left
    \hline
    Colore del jumper & Collegamento \\ [.5ex]
    \hline\hline
    tutti i jumper neri & collegati a GND \\
    rosso & Pin ESP: Vcc a 3.3V \\
    blu & Pin ESP: CH-PD a 3.3V \\
    verde & Pin ESP: RX al Pin Arduino: 10 \\
    arancione & Pin ESP: TX al Pin Arduino: 11 \\
    \hline
\end{tabular} \end{center}
%-----------------------------------------------------------------------

%definizione variabile o funzione
\newcommand \variabileDiProva {55}
\newcommand {\funzioneDiProva}[2] {argomenti = #1 e #2}
%-----------------------------------------------------------------------

%nota a piè pagina
\newcommand{\citAAA}{\nota{Nome} [https://www.link]}
%-----------------------------------------------------------------------

%scrivere testo con _ o %
\verb|porva_p_o_r_v%%%%a|
%-----------------------------------------------------------------------

%piccola riga divisoria
\noindent\makebox[\linewidth]{\rule{15mm}{.4mm}}
%-----------------------------------------------------------------------

%capitoli, sezioni, ecc...
\begin{itemize}[noitemsep]
    \item{}
    \item{}
\end{itemize}
%-----------------------------------------------------------------------

%codice programmazione ARUDINO
\begin{lstlisting}[language=Arduino]

\end{lstlisting}
%divisorio opzionale
\noindent\makebox[\linewidth]{\rule{15mm}{.4mm}}
%-----------------------------------------------------------------------

%codice programmazione PHP
\begin{lstlisting}[language=PHP]
<?php

?>
\end{lstlisting}


%-----------------------------------------------------------------------

%esempio

%-----------------------------------------------------------------------

%esempio

%-----------------------------------------------------------------------
